% preamble.tex --- shared setup for the maya book.
% Targets pdflatex (no fontspec / unicode-math). Safe on every TeX Live.

\usepackage[utf8]{inputenc}
\usepackage[T1]{fontenc}
\usepackage{lmodern}
\usepackage{microtype}

\usepackage[a4paper,margin=1in,headheight=14pt]{geometry}
\usepackage{xcolor}
\usepackage{graphicx}
\usepackage{amsmath,amssymb,amsthm}
\usepackage{enumitem}
\usepackage{hyperref}
\usepackage{fancyhdr}
\usepackage{titlesec}
\usepackage{tcolorbox}
\usepackage{listings}
\usepackage{caption}

\tcbuselibrary{breakable,skins}

% ---- Colour palette (terminal-themed) -----------------------------------
\definecolor{mayablue}{HTML}{2E5C8A}
\definecolor{mayalavender}{HTML}{6E5A8A}
\definecolor{mayagreen}{HTML}{3F7B5A}
\definecolor{mayarose}{HTML}{8A3F5A}
\definecolor{mayaink}{HTML}{1A1A22}
\definecolor{mayagrey}{HTML}{5E5E68}
\definecolor{mayalight}{HTML}{F2F2F4}
\definecolor{mayaamber}{HTML}{B07A1A}

% ---- Hyperref -----------------------------------------------------------
\hypersetup{
  colorlinks=true,
  linkcolor=mayablue,
  citecolor=mayablue,
  urlcolor=mayalavender,
  pdftitle={The Maya Book},
  pdfauthor={agentty docs},
  pdfsubject={A complete tour of the maya TUI library},
  bookmarksnumbered=true,
  bookmarksopen=true
}

% ---- Headers / footers --------------------------------------------------
\pagestyle{fancy}
\fancyhf{}
\fancyhead[LE]{\small\textit{\nouppercase{\leftmark}}}
\fancyhead[RO]{\small\textit{\nouppercase{\rightmark}}}
\fancyfoot[C]{\thepage}
\renewcommand{\headrulewidth}{0.4pt}
\renewcommand{\footrulewidth}{0pt}

% ---- Chapter / section titles ------------------------------------------
\titleformat{\chapter}[display]
  {\normalfont\sffamily\Huge\bfseries\color{mayablue}}
  {\filleft\Large\textit{Chapter \thechapter}}
  {1ex}
  {\titlerule\vspace{1ex}\filright}
\titlespacing*{\chapter}{0pt}{-30pt}{20pt}

\titleformat{\section}{\normalfont\sffamily\Large\bfseries\color{mayalavender}}{\thesection}{1em}{}
\titleformat{\subsection}{\normalfont\sffamily\large\bfseries\color{mayagreen}}{\thesubsection}{1em}{}

% ---- Listings: a C++ flavour we control ---------------------------------
% pdflatex's listings package can't render arbitrary Unicode; we keep code
% strictly ASCII inside lstlistings. For pretty terminal glyphs we drop
% into tcolorbox bodies.
\lstdefinelanguage{maya-cpp}{
  language=C++,
  morekeywords={
    constexpr,consteval,concept,requires,co_await,co_yield,co_return,
    char32_t,char8_t,string_view,expected,variant,visit,move_only_function,
    Signal,Computed,Effect,Batch,Cmd,Sub,Element,Model,Msg,Program,
    BoxNode,TextNode,DynNode,RuntimeTextNode,MapNode,Style,CTStyle,
    Color,Border,BorderStyle,FlexDirection,Align,Justify,Ctx,RunConfig,
    Canvas,StylePool,LayoutNode,Event,Key,Mouse,Resize,Paste,
    overload,Str,FixedName,Refined,EffectSet,Profile,Decision
  }
}
\lstdefinestyle{maya}{
  language=maya-cpp,
  basicstyle=\ttfamily\footnotesize,
  keywordstyle=\color{mayablue}\bfseries,
  stringstyle=\color{mayarose},
  commentstyle=\color{mayagrey}\itshape,
  showstringspaces=false,
  breaklines=true,
  breakatwhitespace=false,
  tabsize=2,
  columns=fullflexible,
  keepspaces=true,
  upquote=true,
  numbers=left,
  numberstyle=\ttfamily\scriptsize\color{mayaink},
  numbersep=12pt,
  frame=leftline,
  framerule=2pt,
  rulecolor=\color{mayablue},
  xleftmargin=3em,
  framexleftmargin=2.5em,
  framesep=8pt,
  literate=%
    {->}{{$\rightarrow$}}2
    {=>}{{$\Rightarrow$}}2
}
\lstset{style=maya}

\lstdefinestyle{shell}{
  language=bash,
  basicstyle=\ttfamily\footnotesize,
  keywordstyle=\color{mayablue},
  commentstyle=\color{mayagrey}\itshape,
  showstringspaces=false,
  breaklines=true,
  frame=leftline,
  framerule=2pt,
  rulecolor=\color{mayagreen},
  xleftmargin=1.5em,
  framesep=8pt,
  numbers=none
}

% ---- Boxes: idea, warning, terminal, theory ----------------------------
\newtcolorbox{idea}[1][]{
  enhanced, breakable,
  colback=mayablue!5, colframe=mayablue, fonttitle=\bfseries,
  title={Idea}, #1
}
\newtcolorbox{theory}[1][]{
  enhanced, breakable,
  colback=mayalavender!7, colframe=mayalavender, fonttitle=\bfseries,
  title={Theory}, #1
}
\newtcolorbox{warning}[1][]{
  enhanced, breakable,
  colback=mayaamber!10, colframe=mayaamber, fonttitle=\bfseries,
  title={Pitfall}, #1
}
\newtcolorbox{recap}[1][]{
  enhanced, breakable,
  colback=mayagreen!8, colframe=mayagreen, fonttitle=\bfseries,
  title={Recap}, #1
}
\newtcolorbox{terminalbox}[1][]{
  enhanced, breakable,
  colback=mayaink, colframe=mayaink, coltext=mayalight,
  fonttitle=\bfseries\color{mayalight},
  title={Terminal}, fontupper=\ttfamily\small, #1
}

% ---- Misc convenience ---------------------------------------------------
\newcommand{\maya}{\textsf{maya}}
\newcommand{\agentty}{\textsf{agentty}}
\newcommand{\cpp}{C\nolinebreak\hspace{-.05em}\raisebox{.2ex}{\small{++}}}
\newcommand{\code}[1]{\texttt{#1}}
\newcommand{\file}[1]{\textsf{\textit{#1}}}
\newcommand{\TODO}[1]{\textcolor{mayarose}{\textbf{TODO:} #1}}

\setlength{\parskip}{0.5\baselineskip}
\setlength{\parindent}{0pt}
